\section{Additive functors}
\label{section-functors}

\noindent
First a completely silly lemma characterizing additive functors
between additive categories.

\begin{lemma}
\label{lemma-additive-functor}
Let $\mathcal{A}$ and $\mathcal{B}$ be additive categories.
Let $F : \mathcal{A} \to \mathcal{B}$ be a functor.
The following are equivalent
\begin{enumerate}
\item $F$ is additive,
\item $F(A) \oplus F(B) \to F(A \oplus B)$ is an isomorphism for
all $A, B \in \mathcal{A}$, and
\item $F(A \oplus B) \to F(A) \oplus F(B)$  is an isomorphism for
all $A, B \in \mathcal{A}$.
\end{enumerate}
\end{lemma}

\begin{proof}
Additive functors commute with direct sums by
Lemma \ref{lemma-additive-additive} hence (1)
implies (2) and (3). If (2) holds and $0$ denotes the
zero object of $\mathcal{A}$, then the map
$F(0) \oplus F(0) \to F(0 \oplus 0) \cong F(0)$
being an isomorphism implies that the diagonal map
$\Hom(F(0), C) \to \Hom(F(0), C) \oplus \Hom(F(0), C)$
is bijective for all $C$ in $\mathcal{B}$. Thus $F(0)$
is the zero object of $\mathcal{B}$ and $F$ maps the
zero morphism between any two objects to the zero morphism.
Similarly holds if (3) holds. It follows that (2) and (3)
are equivalent because the composition
$F(A) \oplus F(B) \to F(A \oplus B) \to F(A) \oplus F(B)$
is given by the matrix
$$
\left(
\begin{matrix}
F(1) & F(0) \\
F(0) & F(1)
\end{matrix}
\right)
$$
and therefore is the identity map. Assume (2) and (3) hold.
Let $f, g : A \to B$ be maps. Then $f + g$ is equal to
the composition
$$
A \to A \oplus A \xrightarrow{\text{diag}(f, g)} B \oplus B \to B
$$
Apply the functor $F$ and consider the following diagram
$$
\xymatrix{
F(A) \ar[r] \ar[rd] &
F(A \oplus A) \ar[rr]_{F(\text{diag}(f, g))} & &
F(B \oplus B) \ar[r] \ar[d] &
F(B) \\
&
F(A) \oplus F(A) \ar[u] \ar[rr]^{\text{diag}(F(f), F(g))} & &
F(B) \oplus F(B) \ar[ru]
}
$$
We claim this is commutative. For the middle square we can verify it
separately for each of the four induced maps $F(A) \to F(B)$
where it follows from the fact that $F$ is a functor and
transforms zero morphisms into zero morphisms.
For the triangle on the left, we use that
$F(A \oplus A) \to F(A) \oplus F(A)$ is an isomorphism
to see that it suffice to check after composition with
this map and this check is trivial. Dually for the other triangle.
Thus going around the bottom is equal to $F(f + g)$ and we conclude.
\end{proof}

\noindent
Recall that we defined, in
Categories, Definition \ref{categories-definition-exact}
the notion of a ``right exact'', ``left exact'' and
``exact'' functor in the setting of a functor between
categories that have finite (co)limits. Thus this
applies in particular to functors between abelian
categories.

\begin{lemma}
\label{lemma-exact-functor}
Let $\mathcal{A}$ and $\mathcal{B}$ be abelian categories.
Let $F : \mathcal{A} \to \mathcal{B}$ be a functor.
\begin{enumerate}
\item If $F$ is either left or right exact, then it is additive.
\item $F$ is left exact if and only if
for every short exact sequence
$0 \to A \to B \to C \to 0$
the sequence $0 \to F(A) \to F(B) \to F(C)$
is exact.
\item $F$ is right exact if and only if for every short exact sequence
$0 \to A \to B \to C \to 0$
the sequence $F(A) \to F(B) \to F(C) \to 0$
is exact.
\item $F$ is exact if and only if for every short exact sequence
$0 \to A \to B \to C \to 0$
the sequence $0 \to F(A) \to F(B) \to F(C) \to 0$
is exact.
\end{enumerate}
\end{lemma}

\begin{proof}
If $F$ is left exact, i.e., $F$ commutes with finite limits, then
$F$ sends products to products, hence $F$ preserved direct sums,
hence $F$ is additive by Lemma \ref{lemma-additive-functor}.
On the other hand, suppose that for every short exact sequence
$0 \to A \to B \to C \to 0$ the sequence $0 \to F(A) \to F(B) \to F(C)$
is exact. Let $A, B$ be two objects. Then we have a short
exact sequence
$$
0 \to A \to A \oplus B \to B \to 0
$$
see for example Lemma \ref{lemma-additive-cat-biproduct-kernel}.
By assumption, the lower row in the commutative diagram
$$
\xymatrix{
0 \ar[r] &
F(A) \ar[d] \ar[r] &
F(A) \oplus F(B) \ar[r] \ar[d] &
F(B) \ar[d] \ar[r] &
0 \\
0 \ar[r] &
F(A) \ar[r] &
F(A \oplus B) \ar[r] &
F(B)
}
$$
is exact. Hence by the snake lemma (Lemma \ref{lemma-snake})
we conclude that $F(A) \oplus F(B) \to F(A \oplus B)$ is an
isomorphism. Hence $F$ is additive in this case as well.
Thus for the rest of the proof we may assume $F$ is additive.

\medskip\noindent
Denote $f : B \to C$ a map from $B$ to $C$.
Exactness of $0 \to A \to B \to C$ just means that
$A = \Ker(f)$. Clearly the kernel of $f$ is
the equalizer of the two maps $f$ and $0$ from $B$ to $C$.
Hence if $F$ commutes with limits, then $F(\Ker(f))
= \Ker(F(f))$ which exactly means that
$0 \to F(A) \to F(B) \to F(C)$ is exact.

\medskip\noindent
Conversely, suppose that $F$ is additive and
transforms any short exact sequence $0 \to A \to B \to C \to 0$ into
an exact sequence $0 \to F(A) \to F(B) \to F(C)$.
Because it is additive it commutes with direct sums
and hence finite products in $\mathcal{A}$. To show
it commutes with finite limits it therefore
suffices to show that it commutes with
equalizers. But equalizers in an abelian category
are the same as the kernel of the difference map,
hence it suffices to show that $F$ commutes with
taking kernels. Let $f : A \to B$ be a morphism.
Factor $f$ as $A \to I \to B$ with $f' : A \to I$ surjective
and $i : I \to B$ injective. (This is possible by the
definition of an abelian category.) Then it is
clear that $\Ker(f) = \Ker(f')$. Also
$0 \to \Ker(f') \to A \to I \to 0$
and
$0 \to I \to B \to B/I \to 0$
are short exact. By the condition imposed on $F$
we see that
$0 \to F(\Ker(f')) \to F(A) \to F(I)$
and
$0 \to F(I) \to F(B) \to F(B/I)$
are exact. Hence it is also the case that
$F(\Ker(f'))$ is the kernel of the map
$F(A) \to F(B)$, and we win.

\medskip\noindent
The proof of (3) is similar to the proof of (2).
Statement (4) is a combination of (2) and (3).
\end{proof}

\begin{lemma}
\label{lemma-exact-functor-ext}
Let $\mathcal{A}$ and $\mathcal{B}$ be abelian categories.
Let $F : \mathcal{A} \to \mathcal{B}$ be an exact functor.
For every pair of objects $A, B$ of $\mathcal{A}$ the
functor $F$ induces an abelian group homomorphism
$$
\Ext_\mathcal{A}(B, A)
\longrightarrow
\Ext_\mathcal{B}(F(B), F(A))
$$
which maps the extension $E$ to $F(E)$.
\end{lemma}

\begin{proof}
Omitted.
\end{proof}

\noindent
The following lemma is used in the proof that the category of abelian
sheaves on a site is abelian, where the functor $b$ is sheafification.

\begin{lemma}
\label{lemma-adjoint-get-abelian}
Let $a : \mathcal{A} \to \mathcal{B}$ and $b : \mathcal{B} \to \mathcal{A}$
be functors. Assume that
\begin{enumerate}
\item $\mathcal{A}$, $\mathcal{B}$ are additive categories,
$a$, $b$ are additive functors, and $a$ is right adjoint to $b$,
\item $\mathcal{B}$ is abelian and $b$ is left exact, and
\item $ba \cong \text{id}_\mathcal{A}$.
\end{enumerate}
Then $\mathcal{A}$ is abelian.
\end{lemma}

\begin{proof}
As $\mathcal{B}$ is abelian we see that all finite limits and colimits
exist in $\mathcal{B}$ by Lemma \ref{lemma-colimit-abelian-category}.
Since $b$ is a left adjoint we see that $b$ is also right exact
and hence exact, see
Categories, Lemma \ref{categories-lemma-exact-adjoint}.
Let $\varphi : B_1 \to B_2$ be a morphism of $\mathcal{B}$.
In particular, if $K = \Ker(B_1 \to B_2)$, then $K$ is
the equalizer of $0$ and $\varphi$ and hence
$bK$ is the equalizer of $0$ and $b\varphi$, hence
$bK$ is the kernel of $b\varphi$. Similarly, if
$Q = \Coker(B_1 \to B_2)$, then $Q$ is
the coequalizer of $0$ and $\varphi$ and hence
$bQ$ is the coequalizer of $0$ and $b\varphi$, hence
$bQ$ is the cokernel of $b\varphi$. Thus we see that every morphism
of the form $b\varphi$ in $\mathcal{A}$ has a kernel and a cokernel.
However, since $ba \cong \text{id}$ we see that every morphism of
$\mathcal{A}$ is of this form, and we conclude that kernels and
cokernels exist in $\mathcal{A}$. In fact, the argument shows that
if $\psi : A_1 \to A_2$ is a morphism then
$$
\Ker(\psi) = b\Ker(a\psi),
\quad\text{and}\quad
\Coker(\psi) = b\Coker(a\psi).
$$
Now we still have to show that $\Coim(\psi)= \Im(\psi)$.
We do this as follows.
First note that since $\mathcal{A}$ has kernels and cokernels it
has all finite limits and colimits (see proof of
Lemma \ref{lemma-colimit-abelian-category}).
Hence we see by Categories, Lemma \ref{categories-lemma-exact-adjoint}
that $a$ is left exact and
hence transforms kernels (=equalizers) into kernels.
\begin{align*}
\Coim(\psi)
& =
\Coker(\Ker(\psi) \to A_1)
& \text{by definition} \\
& =
b\Coker(a(\Ker(\psi) \to A_1))
& \text{by formula above} \\
& =
b\Coker(\Ker(a\psi) \to aA_1))
& a\text{ preserves kernels} \\
& =
b\Coim(a\psi)
& \text{by definition} \\
& =
b\Im(a\psi)
& \mathcal{B}\text{ is abelian} \\
& =
b\Ker(aA_2 \to \Coker(a\psi))
& \text{by definition} \\
& =
\Ker(baA_2 \to b\Coker(a\psi))
& b\text{ preserves kernels} \\
& =
\Ker(A_2 \to b\Coker(a\psi))
& ba = \text{id}_\mathcal{A} \\
& =
\Ker(A_2 \to \Coker(\psi))
& \text{by formula above} \\
& =
\Im(\psi)
& \text{by definition}
\end{align*}
Thus the lemma holds.
\end{proof}




